%!TEX TS-program = pdflatex
%!TEX encoding = UTF-8 Unicode
%!TEX spellcheck = English (Aspell)

\documentclass[11pt]{article}
\usepackage{lucy,handout,graphicx}

% =========================================================
\begin{document}


\noindent 
A robot can operate in one of $n$ modes. You are given a list of $m$ commands that change
the robot from one mode to another; the reverse command may not exist.
The
manufacturer claims: ``No matter which commands you execute, you can always
restore the robot to the mode it was in before you started.'' 
Give an efficient algorithm to check this claim. (Note: there is an example input
at the bottom of the page, but you should answer these questions for a general input.)

\begin{itemize}\itemsep40pt
\item What are the vertices?  What does each vertex represent?
\item What are the edges?  Are they directed or undirected?
\item If the vertices and/or edges have associated values, what are they?
\item \textbf{What is the precise relationship between your graph and the stated problem?}
\item What standard graph algorithm are you using to solve the problem?
\item What is the running time of your entire algorithm, \emph{including} the time to build the graph, \emph{as a function of the original input parameters}?
\end{itemize}

\vspace{1in}


Here is an example you may want to explore:

\begin{center}
\begin{tabular}{lll}
\hline
Command & Starting mode & Resulting mode \\
\hline
Begin charging  & Standby   & Charging  \\
Stop charging   & Charging  & Standby   \\
Switch to mop   & Vacuuming & Mopping   \\
Begin drying    & Mopping   & Drying    \\
Return to base  & Drying    & Returning \\
Resume cleaning & Returning & Vacuuming \\
\hline
\end{tabular}
\end{center}





%----------------------------------------------------------------------


\end{document}